\documentclass[11pt,a4paper]{amsart}
\usepackage[utf8]{inputenc}
\usepackage[T1]{fontenc}
\usepackage{amsmath,amssymb,amsthm}
\usepackage{hyperref}

\newtheorem{theorem}{Theorem}[section]
\newtheorem{proposition}[theorem]{Proposition}
\newtheorem{remark}[theorem]{Remark}
\newtheorem{openproblem}[theorem]{Open Problem}

\title{P01 --- BC Prime-Power Weights and Local Arithmetic Structure}
\author{Objekt-X-Programm}
\date{2026-08-09}

\begin{document}
\maketitle

\begin{abstract}
We record the reconciled local Bost--Connes prime-power layer used as a
prerequisite for P09.  The primitive $p$-channel produces the algebraic factor
$\log p/\sqrt p$, and the arithmetic identity
$\Lambda(p^m)/\sqrt{p^m}=\log p/p^{m/2}$ is unconditional.  The stronger
all-$n$ operator realization and the Hilbert-space functional calculus remain
open.  The support-separation statement from NEU-250j is retained without
upgrading it to an orthogonality theorem.
\end{abstract}

\section{Primitive BC prime channel}

\begin{proposition}[Primitive local factor]
On the explicitly treated primitive channel $j_R=e_RV_p$, the audited
algebraic relations are
\[
  h_p^{\rm bal}=p^{-1/2}I,
  \qquad
  H_{\rm BC}j_R=(\log p)j_R.
\]
Hence the primitive local contribution carries the factor
\[
  \boxed{\frac{\log p}{\sqrt p}}.
\]
This statement is only an algebraic/partial realization.  Self-adjointness,
closure, domains and a global Hilbert functional calculus for $H_{\rm BC}$
are not supplied here.
\end{proposition}

\section{Prime-power arithmetic}

\begin{theorem}[Arithmetic prime-power identity]
For every prime $p$ and every integer $m\ge1$,
\[
  \boxed{
  \frac{\Lambda(p^m)}{\sqrt{p^m}}
  =\frac{\log p}{p^{m/2}}.
  }
\]
This is a purely arithmetic and RH-free identity.
\end{theorem}

\begin{openproblem}[All-$n$ operator realization]
The general operator formula
\[
  h_n^{\rm bal}=n^{-1/2}I\qquad(n\ge1)
\]
is not proved in the audited source cone.  NEU-250g computes the primitive
$p$-channel, whereas the all-$n$ extension invoked in NEU-250i lacks an
independent general lemma.  Consequently the full prime-power operator
realization through $h_{p^m}^{\rm bal}$ or $H_{\rm pr}$ remains conditional.
No global operator identity $H_{\rm pr}=\Lambda$ is asserted.
\end{openproblem}

\section{Support separation}

\begin{theorem}[Mangoldt support versus direct cross-prime collision]
Let $p\ne q$ be primes and suppose
\[
  pm_p=qm_q=M.
\]
Then $M$ has at least two distinct prime divisors and therefore
\[
  \Lambda(M)=0.
\]
Thus, for the direct cross-prime collision considered in NEU-250j,
\[
  \boxed{
  \operatorname{supp}\Lambda
  \cap
  \operatorname{supp}(\text{direct cross-prime collision})
  =\varnothing.
  }
\]
\end{theorem}

\begin{remark}[Firewall]
Support separation is not a theorem of Hilbert-space orthogonality or of
pairwise disjoint prime-channel images.  The nonorthogonal prime-channel
geometry retained in P05 is unaffected.
\end{remark}

\section{Status and provenance}

The authoritative reconciliation is
\texttt{audits/AUDIT-2026-08-09\_P01\_Dependency\_Reaudit\_vor\_P09.md},
based on the doubly checked F4 end state of the NEU-250f--r chain.
The canonical Markdown source is
\texttt{papers/P01\_BC\_Prime\_Power\_Weights.md}.

\[
  \boxed{\text{P01 --- SYN FINAL AUDITED; dependency-reconciled for P09}.}
\]

No Riemann-hypothesis proof, Object-X construction, or global
Hilbert--P\'olya realization is claimed.

\end{document}
